\documentclass{article}
\usepackage[UTF8]{ctex}
\usepackage{amsmath,amssymb,amsfonts}
\usepackage{enumitem}
\usepackage{geometry}
\geometry{a4paper, margin=1in}
% === 错题本排版优化（不参与题目标识） ===
\setlength{\parindent}{0pt}
\setlength{\parskip}{0.35em plus 0.10em minus 0.08em}
\setlength{\emergencystretch}{3em}
\setlist{itemsep=0.25em, topsep=0.40em, parsep=0pt}
\allowdisplaybreaks[2]
% === 错题本排版优化结束 ===
\begin{document}
\section*{广西大学2025年高等代数真题}
\section*{153．广西大学2025年高等代数试卷（回忆版）}
一．填空题．每题6分，共48分．\\

\subsection{题目1}
1．设 \(A, B\) 为 \(n \times n\) 矩阵，\(|A|=2,|B|=-3\) ，则 \(\left|2 A^{*} B^{-1}\right|=\) \(\_\_\_\_\)．\\
\(\_\_\_\_\)

\subsection{题目2}
2．设向量组 \(\alpha_{1}=(3,-2,0)^{\mathrm{T}}, \alpha_{2}=(27,-18,0)^{\mathrm{T}}, \alpha_{3}=(-1,5,8)^{\mathrm{T}}\)（注．另一个版本 \(\alpha_{2}\) 的第二个分量是18），秩为 \(\_\_\_\_\)，极大线性无关组为 \(\_\_\_\_\)．

\subsection{题目3}
3．设 \(A=\left(\begin{array}{llll}0 & 0 & 1 & 2 \\ 0 & 0 & 1 & 3 \\ 2 & 1 & 0 & 0 \\ 5 & 3 & 0 & 0\end{array}\right)\) ，则 \(A^{-1}=\)

\subsection{题目4}
4．设二次型 \(f\left(x_{1}, x_{2}, x_{3}\right)=x_{1}^{2}+k x_{2}^{2}+(k-2) x_{3}^{2}+2 x_{1} x_{2}\) 正定，求 \(k\) 的取值范围 \(\_\_\_\_\)．

\subsection{题目5}
5．设 \(f(x)=x^{4}-5 x^{3}+a x^{2}+b x+9 \in \mathbb{Q}[x]\) ，若3是 \(f(x)\) 的二重根，则 \(a=\) \(\_\_\_\_\)，\(b=\) \(\_\_\_\_\)．\\
\(\_\_\_\_\)

\subsection{题目6}
6．设 \(V\) 为数域 \(K\) 上的线性空间，\(V_{1}, V_{2}\) 为 \(V\) 的子空间， \(\operatorname{dim} V=9, \operatorname{dim} V_{1}=5, \operatorname{dim} V_{2}=6\) ，则 \(\operatorname{dim}\left(V_{1} \cap V_{2}\right)\) 的最小值为 \(\_\_\_\_\)．

\begin{itemize}

\subsection{题目7}
\end{itemize}

\subsection{题目7}
原题边界“7．”：设线性变换 \(\sigma\) 在 \(V\) 的一组基 \(\alpha_{1}, \alpha_{2}, \alpha_{3}\) 下的矩阵是
\[
A=\left(\begin{array}{ccc}
2 & -2 & 2 \\
-2 & -1 & 4 \\
2 & 4 & 1
\end{array}\right)
\]
求 \(\sigma\) 的全部特征值 \(\_\_\_\_\)、

\subsection{题目8}

\subsection{题目8}
原题边界“8．”：欧氏空间 \(V\) 的基 \(\alpha_{1}, \alpha_{2}, \alpha_{3}\) 的度量矩阵是
\[
\left(\begin{array}{ccc}
2 & 0 & -1 \\
0 & 4 & 0 \\
-1 & 0 & 1
\end{array}\right)
\]
\(\beta=2 \alpha_{1}-\alpha_{2}+\alpha_{3}\) ，求 \(|\beta|=\) \(\_\_\_\_\)．

二．（15 分）求下列 \(n\) 阶行列式

\[
D_{n}=\left|\begin{array}{ccccc}
1 & 1 & 1 & \cdots & 1 \\
a_{1} & a_{2} & a_{3} & \cdots & a_{n} \\
a_{1}^{2} & a_{2}^{2} & a_{3}^{2} & \cdots & a_{n}^{2} \\
\vdots & \vdots & \vdots & & \vdots \\
a_{1}^{n-1} & a_{2}^{n-1} & a_{3}^{n-1} & \cdots & a_{n}^{n-1}
\end{array}\right|(n \geq 2) .
\]

三。（15 分）设 \(A\) 是数域 \(K\) 上的 \(n \times n\) 矩阵，\(A\) 可对角化，证明：\(A\) 的复特征值都属于 \(K\) ，且属于不同特征值的对应特征子空间的维数等于每个特征值的重数。（注。另一个版本是证明矩阵的每个特征值的几何重数小于等于代数重数）

四．（15 分）设 \(f(x)=a_{n} x^{n}+a_{n-1} x^{n-1}+\cdots+a_{1} x+a_{0}\) 是一个 \(n\) 次整系数多项式，且存在素数 \(p\) ，使得

\[
p \nmid a_{n} ; p \mid a_{n-1}, \cdots, a_{1}, a_{0}, p^{2} \nmid a_{0} .
\]

则 \(f(x)\) 在有理数域上不可约．

五．（15 分）\(A, B\) 是 \(n\) 级实对称矩阵，\(A B=B A\) ，证明：存在正交矩阵 \(T\) ，使得 \(T^{\mathrm{T}} A T, T^{\mathrm{T}} B T\) 同时为对角阵。

六。（16 分）\(\sigma\) 是线性空间 \(V\) 上的线性变换，其在 \(V\) 的一组基 \(\alpha_{1}, \alpha_{2}, \alpha_{3}\) 下的矩阵为

\[
A=\left(\begin{array}{lll}
1 & 0 & 0 \\
1 & 2 & 1 \\
1 & 0 & 1
\end{array}\right) . \text { (注. 另一版本的左下角元素为 }-1 \text { ) }
\]

\begin{itemize}
  \item[（1）] 求 \(\sigma\) 的最小多项式 \(m(\lambda)\) ．
  \item[（2）] 对应于 \(m(\lambda)\) 的因式分解，写出 \(V\) 的直和分解式，并求出分解式中出现的每一个子空间的基。
\end{itemize}

七．（16 分）\(\sigma\) 是欧氏空间 \(V\) 上的对称变换，且 \(\sigma^{2}=\sigma\) 。证明：\(\sigma\) 是从 \(V\) 到 \(\sigma\) 像空间 \(\operatorname{Im} \sigma\) 的正交投影。

八．（10 分）\(A\) 为 \(n\) 级矩阵，\(E\) 为 \(n\) 级单位矩阵，\(A^{4}=E\) ，证明：

\[
r(E-A)+r\left(E+A+A^{2}+A^{3}\right)=n .
\]
\end{document}
